\chapter{Maryam Mirzakhani}\label{ch2014maryammirzakhani:maryam-mirzakhani-fields-medal-2014-a-comprehensive-report-on-her-work-and-impact}

\begin{center}\rule{0.5\linewidth}{0.5pt}\end{center}

\textbf{Prepared:} July 2026 \textbf{Subject:} Maryam Mirzakhani (May 12, 1977 -- July 14, 2017, Tehran, Iran) \textbf{Position:} Professor of Mathematics, Stanford University (2008--2017) \textbf{Doctoral Advisor:} Curtis T. McMullen (Fields Medalist, 1998) \textbf{Fields Medal Citation:} \emph{``For outstanding contributions to the dynamics and geometry of Riemann surfaces and their moduli spaces.''}

\begin{center}\rule{0.5\linewidth}{0.5pt}\end{center}

\begin{prereqs} Hyperbolic geometry (geodesics, Fenchel-Nielsen coordinates), basic Teichm\"uller theory, ergodic theory, algebraic geometry (moduli spaces). \end{prereqs}

\begin{objectives} Explain Mirzakhani's recurrence \cite{mirzakhani2007} theorem for the geodesic flow; describe the counting of simple closed geodesics; understand the volume recursion for moduli spaces and the Magic Wand Theorem. \end{objectives}

\section{Problem History}

Riemann surfaces of genus $g \ge 2$ admit hyperbolic metrics of constant curvature $-1$ by the Uniformization Theorem: any such surface can be realized as a quotient $\mathbb{H}^2/\Gamma$ where $\mathbb{H}^2$ is the hyperbolic plane and $\Gamma$ is a discrete group of isometries. The moduli space $\mathcal{M}_g$ parametrizes all such surfaces up to biholomorphism, and its compactification $\overline{\mathcal{M}}_g$ (the Deligne--Mumford compactification) includes surfaces with nodal singularities. The dimension of $\mathcal{M}_g$ is $6g-6$.

A fundamental question in hyperbolic geometry concerned the growth rate of \textbf{simple closed geodesics}---geodesics that form closed loops without self-intersection. The ``prime number theorem for geodesics'' (Margulis, 1960s) shows that the number of \emph{all} closed geodesics of length $\le L$ grows exponentially: $\pi(L) \sim e^{hL}/(hL)$ where $h$ is the topological entropy. But for \emph{simple} closed geodesics, the counting problem was much more difficult because the simplicity constraint makes the counting highly nontrivial.

The Weil--Petersson metric on $\mathcal{M}_g$ provides a natural Riemannian structure with a finite volume form. Computing the volume $V_g$ was a classical problem solved only for small genus: $V_2 = \pi^4/270$, $V_3 = 11\pi^6/37800$. For general genus, the problem appeared intractable.

The Weil--Petersson metric is defined using the $L^2$ inner product on quadratic differentials: for tangent vectors $\mu, \nu \in T_X \mathcal{M}_g$ (which are Beltrami differentials on $X$), the metric is $\langle \mu, \nu \rangle_{WP} = \int_X \mu \overline{\nu} \rho^{-2}$, where $\rho$ is the hyperbolic metric on $X$. The volume form is $(WP)^{3g-3}$, and the total volume $V_g = \int_{\mathcal{M}_g} (WP)^{3g-3}$ is a rational multiple of $\pi^{6g-6}$. The Weil--Petersson metric has negative sectional curvature and is incomplete: geodesics can reach the boundary $\overline{\mathcal{M}}_g \setminus \mathcal{M}_g$ in finite time, corresponding to degenerating Riemann surfaces.

The Weil--Petersson volume has remarkable properties. Wolpert showed that the Weil--Petersson symplectic form $\omega_{WP}$ satisfies $\omega_{WP} = \sum_{i=1}^{3g-3} d\alpha_i \wedge d\beta_i$, where $(\alpha_i, \beta_i)$ are Fenchel--Nielsen coordinates (length and twist parameters for a pants decomposition). This symplectic structure connects the Weil--Petersson metric to the Hamiltonian mechanics of surface deformations. The volume $V_g$ is computed by integrating the top power of the symplectic form over $\mathcal{M}_g$, and the resulting numbers encode deep information about the topology of the moduli space.

The Weil--Petersson metric has negative sectional curvature, which means that geodesics in the moduli space diverge exponentially. This negative curvature is responsible for the ergodic properties of the Teichm\"uller geodesic flow and the earthquake flow. The curvature tensor of the Weil--Petersson metric can be expressed in terms of the geometry of the Riemann surface, providing a deep connection between the metric geometry of the moduli space and the hyperbolic geometry of the surfaces it parametrizes.

The \textbf{Kontsevich--Witten formula} (1992), proved by Maxim Kontsevich for which he received the Fields Medal in 1998, relates the intersection theory of moduli spaces to the KdV hierarchy of integrable PDEs. The intersection numbers $\langle\tau_{d_1}\cdots\tau_{d_n}\rangle_g = \int_{\overline{\mathcal{M}}_{g,n}}\psi_1^{d_1}\cdots\psi_n^{d_n}$ satisfy the KdV equations. Kontsevich's proof used matrix models and combinatorial methods rather than direct geometry of moduli spaces.

In dynamics, the Teichm\"uller geodesic flow on Teichm\"uller space and the earthquake flow (deforming surfaces by cutting along geodesics and twisting) were known to be ergodic on the Teichm\"uller space. The action of $\mathrm{SL}(2,\mathbb{R})$ on the moduli space of Abelian differentials provides a framework for studying billiard dynamics. Understanding the orbit closures of this action---analogous to understanding invariant submanifolds in symplectic geometry---was a major open problem.

The Teichm\"uller space $\mathcal{T}_g$ is a complex manifold of dimension $3g-3$ parametrizing marked Riemann surfaces. The mapping class group $\mathrm{Mod}_g$ acts properly discontinuously on $\mathcal{T}_g$, and the quotient is the moduli space $\mathcal{M}_g$. The Teichm\"uller geodesic flow is a one-parameter flow on the unit tangent bundle of $\mathcal{T}_g$ that corresponds to stretching along the horizontal direction and contracting along the vertical direction of the quadratic differential. Understanding the ergodic properties of this flow---its mixing rates, invariant measures, and orbit closures---was a major problem in dynamics.

The earthquake flow is another important flow on Teichm\"uller space, introduced by Thurston. It deforms a hyperbolic surface by cutting along a measured lamination and re-gluing with a twist. Unlike the Teichm\"uller flow, the earthquake flow is not defined by a simple geometric operation on differentials, making its analysis much more difficult. Thurston conjectured that the earthquake flow is ergodic on Teichm\"uller space, a result that Mirzakhani later proved using methods from ergodic theory and hyperbolic geometry.

The \textbf{moduli space of Abelian differentials} $\mathcal{H}_g$ parametrizes pairs $(X, \omega)$ where $X$ is a Riemann surface of genus $g$ and $\omega$ is a holomorphic 1-form on $X$. The $\mathrm{SL}(2,\mathbb{R})$ action on $\mathcal{H}_g$ is defined by the real-linear changes of basis for the real and imaginary parts of $\omega$: for $A = \begin{pmatrix} a & b \\ c & d \end{pmatrix} \in \mathrm{SL}(2,\mathbb{R})$, the action is $A \cdot (X, \omega) = (X, (a + bi)\omega)$ when $\omega$ is normalized so that its period lattice is a sublattice of $\mathbb{Z}^2$. The orbit closures of this action correspond to invariant submanifolds of $\mathcal{H}_g$, and the Eskin--Mirzakhani classification shows that all such orbit closures are algebraic subvarieties.

\section{How Was It Solved}

In her PhD thesis (2004), Mirzakhani proved that the number of simple closed geodesics of length $\le L$ on a hyperbolic surface of genus $g$ is asymptotic to $c_g \cdot L^{6g-6}$, where $6g-6 = \dim\mathcal{M}_g$. This polynomial growth---in sharp contrast to the exponential growth of all geodesics---confirmed a long-standing conjecture.

The key insight was to connect counting geodesics to volumes of moduli spaces via a \textbf{generalized McShane identity}. The original McShane identity for a punctured torus states $\sum_{\gamma} 1/(1+e^{\ell(\gamma)}) = 1/2$. Mirzakhani generalized this to bordered hyperbolic surfaces with geodesic boundary components of lengths $x_1,\ldots,x_n$:
\[
\sum_{\{\gamma_1,\gamma_2\}\in\mathcal{F}_i} D(\ell_{\gamma_1},\ell_{\gamma_2},x_i) + \sum_{j\ne i}\sum_{\gamma\in\mathcal{F}_{i,j}} R(\ell_{\gamma},x_i,x_j) = x_i,
\]
where $D$ and $R$ are explicit elementary functions.

Integrating both sides over moduli space transformed this identity into a \textbf{recursive formula for Weil--Petersson volumes}: $V_{g,n}(L_1,\ldots,L_n)$ is expressed in terms of volumes of moduli spaces with smaller genus or fewer boundaries. The recursion shows that $V_{g,n}$ is a polynomial in $L_1^2,\ldots,L_n^2$ with rational coefficients.

By plugging the volume formula into the expansion $V_{g,n} = \sum \frac{L_1^{2d_1}\cdots L_n^{2d_n}}{2^{\sum d_i}d_1!\cdots d_n!}\langle\tau_{d_1}\cdots\tau_{d_n}\rangle_g$, Mirzakhani derived recursive relations for the intersection numbers equivalent to the KdV equations, giving a new geometric proof of the Kontsevich--Witten formula that avoided matrix integrals.

The \textbf{Magic Wand Theorem} (2013, with Alex Eskin, extended with Amir Mohammadi) states that every orbit closure of the $\mathrm{SL}(2,\mathbb{R})$ action on the moduli space of Abelian differentials is a complex-algebraic subvariety. The proof, spanning over 200 pages, combines ergodic theory, Hodge theory, homogeneous dynamics (Ratner's theorem), and Teichm\"uller theory in a measure-rigidity argument.

Mirzakhani also proved the ergodicity of the earthquake flow on Teichm\"uller space (2008), confirming Thurston's conjecture, and established large-genus asymptotics: $V_g \sim \pi^{-1/2}(2/\pi^2)^g(2g)!$, inaugurating the study of random hyperbolic surfaces.

\section{Impact of This Discovery}

The counting result for simple closed geodesics resolved a central problem in hyperbolic geometry. The polynomial growth rate $L^{6g-6}$ directly reflects the dimension of moduli space, revealing a deep connection between geometry and dynamics.

The volume recursion formula and its connection to the Kontsevich--Witten formula showed that the Kontsevich--Witten result is not an isolated phenomenon but is intimately connected to the hyperbolic geometry of Riemann surfaces. This provided new tools for computing intersection numbers and studying the topology of moduli spaces.

The Magic Wand Theorem is one of the most important results in dynamical systems in the 21st century. It revealed that the seemingly chaotic dynamics of moduli spaces are governed by rigid algebraic structures, opening entirely new directions in the study of $\mathrm{SL}(2,\mathbb{R})$-actions.

The classification of invariant and stationary measures (Eskin--Mirzakhani, 2018, posthumous) has significant consequences: it implies the $\mathrm{SL}(2,\mathbb{R})$ action is rigid in a strong sense, provides a complete description of all ergodic behaviors, and has applications to billiards, translation surfaces, and IETs.

As the first woman to win the Fields Medal, Mirzakhani became a symbol of possibility for women in mathematics. Her story---from Tehran to Harvard to the Fields Medal---inspired countless young people worldwide. Initiatives established in her memory include the Maryam Mirzakhani New Frontiers Prize and the May 12 Initiative celebrating women in mathematics.

The counting of simple closed geodesics has applications to the study of lengths of closed geodesics on hyperbolic surfaces, which are related to the lengths of words in the fundamental group of the surface. The polynomial growth rate $L^{6g-6}$ provides a precise count of the number of geodesics of length at most $L$, which has implications for the study of the distribution of lengths of closed geodesics and the statistics of hyperbolic surfaces.

The \textbf{Kontsevich--Witten formula} relates intersection numbers on $\overline{\mathcal{M}}_{g,n}$ to the KdV hierarchy:
\[
\frac{\partial F_g}{\partial t_n} = \frac{1}{2} \sum_{i+j=n-1} (2i+1)(2j+1) t_i t_j \frac{\partial F_g}{\partial t_{i+j}} + \frac{1}{16} \frac{\partial^2 F_g}{\partial t_{n-1}^2} + \frac{1}{32} \frac{\partial F_{g-1}}{\partial t_{n-1}} + \frac{1}{64} \frac{\partial F_{g-1}}{\partial t_{n-2}},
\]
where $F_g = \sum_{n} \frac{1}{n!} \sum_{d_1+\cdots+d_n=3g-3+n} \langle\tau_{d_1}\cdots\tau_{d_n}\rangle_g t_{d_1}\cdots t_{d_n}$. Mirzakhani's proof of this formula using Weil--Petersson volumes provided a geometric interpretation of the KdV hierarchy, connecting integrable systems to the geometry of Riemann surfaces.

The \textbf{Weil--Petersson volume} $V_{g,n}(L_1,\ldots,L_n)$ is a polynomial in $L_1^2,\ldots,L_n^2$ with rational coefficients. The recursion formula expresses $V_{g,n}$ in terms of volumes with smaller genus or fewer boundaries:
\[
V_{g,n}(L_1,\ldots,L_n) = \frac{1}{2}\sum_{i=1}^n \sum_{\mathcal{I}} V_{g-1,n+1}(L_{\mathcal{I}}, L_i, L_i, L_{\mathcal{I}^c}) + \cdots
\]
where the sum is over partitions of $\{1,\ldots,n\}\setminus\{i\}$ into two subsets $\mathcal{I}$ and $\mathcal{I}^c$. This recursion allows the systematic computation of Weil--Petersson volumes, and the resulting numbers have been tabulated for small genus and marked points.

\section{Modern Views and Improvements}

The study of random hyperbolic surfaces, initiated by Mirzakhani's large-genus asymptotics, has become an active research area. Random surfaces chosen with respect to the Weil--Petersson volume form have typical properties: the shortest geodesic has length $\sim 2\log(2g)$, the Cheeger constant is bounded away from zero, and the spectral gap of the Laplacian is bounded below.

The Magic Wand Theorem has been applied to study the dynamics of billiards in rational polygons, the ergodic theory of IETs, and the geometry of flat surfaces. Extensions to higher-rank actions and other moduli spaces are under investigation.

The connection between Weil--Petersson volumes and intersection theory has led to new results on the asymptotics of intersection numbers on moduli spaces (Mirzakhani--Zograf).

The methods developed by Mirzakhani---combining hyperbolic geometry, Teichm\"uller theory, dynamics, and algebraic geometry---have become a paradigm for modern geometric research, influencing work on mapping class groups, geometric group theory, and low-dimensional topology.

Current research builds on Mirzakhani's work in several directions: the classification of higher-rank actions on moduli spaces, the statistical properties of random hyperbolic surfaces, and the development of new counting methods for translation surfaces. The Eskin--Mirzakhani classification has been extended to many new settings, and the Weil--Petersson volume recursion continues to yield new results about intersection numbers on moduli spaces.

The Weil--Petersson volume recursion has also been applied to compute intersection numbers on moduli spaces that were previously inaccessible. The recursion provides a systematic method for computing intersection numbers of psi-classes and kappa-classes on $\overline{\mathcal{M}}_{g,n}$, and the resulting numbers have been tabulated for small genus and marked points. These computations have revealed unexpected patterns and connections to other areas of mathematics, including integrable systems and representation theory.

The work of Mirzakhani on the dynamics of Teichm\"uller space has inspired new research on the ergodic properties of mapping class group actions. The classification of invariant measures, known as the measure classification theorem, has been extended to higher-rank actions and to other moduli spaces. These extensions have implications for the study of translation surfaces, billiards in rational polygons, and the geometry of flat surfaces.

The classification of $\mathrm{SL}(2,\mathbb{R})$-invariant measures on the moduli space of Abelian differentials, completed posthumously with Eskin, has led to new results on the dynamics of translation surfaces. The classification shows that all invariant measures are algebraic---they are supported on orbit closures of the $\mathrm{SL}(2,\mathbb{R})$ action. This rigidity result has implications for the study of billiards in rational polygons, the distribution of saddle connections, and the geometry of flat surfaces. The ergodic theory of translation surfaces continues to develop using the techniques pioneered by Mirzakhani and Eskin.

The Magic Wand Theorem has been applied to the study of dynamics on moduli spaces of higher-dimensional abelian varieties. The classification of orbit closures for the $\mathrm{SL}(2,\mathbb{R})$ action on the moduli space of Abelian differentials has been extended to other rank-one Lie group actions on moduli spaces, including the action of the semisimple Lie group on the moduli space of Higgs bundles. These extensions have implications for the study of non-abelian Hodge theory and for the geometry of moduli spaces of principal bundles on algebraic curves.

The study of the dynamics of the earthquake flow, which Mirzakhani proved is ergodic, has led to new results on the distribution of earthquakes on Teichm\"uller space. The earthquake flow is a one-parameter flow that deforms a hyperbolic surface by cutting along a measured lamination and re-gluing with a twist. The ergodicity of the earthquake flow implies that typical earthquakes visit every region of Teichm\"uller space, which has implications for the geometry of the moduli space and for the distribution of hyperbolic structures on surfaces.

The study of mapping class group dynamics on Teichm\"uller space has been influenced by Mirzakhani's work. The classification of orbit closures of the $\mathrm{SL}(2,\mathbb{R})$ action on the moduli space of Abelian differentials has inspired analogous questions for other group actions on moduli spaces. The development of new techniques for studying the dynamics of mapping class group actions, building on Mirzakhani's methods, is an active area of research with applications to the geometry of surfaces and to the theory of Teichm\"uller dynamics.

The study of the Thurston norm on the first cohomology of a 3-manifold has been influenced by Mirzakhani's work on the geometry of the mapping class group. The Thurston norm provides a natural polyhedral structure on the first cohomology of a 3-manifold, and its relationship to the dynamics of the mapping class group is an active area of research. The techniques developed by Mirzakhani for studying the geometry of Teichm\"uller space have been applied to the study of the Thurston norm and to the geometry of 3-manifolds.

The study of the Weil-Petersson metric has led to new results on the geometry of the moduli space of curves. The curvature properties of the Weil-Petersson metric, including the pinching of the sectional curvature and the growth of geodesic balls, have been studied using Mirzakhani's volume recursion. The connection between the Weil-Petersson metric and the theory of random hyperbolic surfaces, through the large-genus asymptotics, is an active area of research that builds on Mirzakhani's foundational work.

\begin{center}\rule{0.5\linewidth}{0.5pt}\end{center}

\section*{Summary}
\addcontentsline{toc}{section}{Summary}

Mirzakhani proved the polynomial growth of simple closed geodesics via a volume recursion for moduli spaces, gave a new proof of the Kontsevich--Witten formula, and established the Magic Wand Theorem classifying orbit closures of the $\mathrm{SL}(2,\mathbb{R})$ action.
